\section{Dos umbrales de longitud y una sola ley de historias}
Sea $M=\lceil\delta n\rceil$. La ley es $\nu$ con orden de filas uniforme revelado
progresivamente. Una salida se cuenta en el hijo, incluso si coincide con el suelo.
Tomamos una vez la clase de C, \S5.4. Sus constantes incluyen $C,\gamma>0$
para el punto de reinas $p$, con $\max p_e\le C/j$ y brecha de filas y columnas
$\ge\gamma$ en el interior. Su umbral geométrico se denota $\alpha_{\rm hist}$.
No se cambia el suceso histórico de C al introducir otra longitud $\alpha j$ en el extremo.
El Teorema 8.1 de B da predictor finito para todos los hijos legales del padre regular.
En su prueba, $e_c=L(B+c)-L(B)-\log p_c$, $\delta_c=O(1/j)$ y
$e_c-\log(1+\delta_c)=O(j^{-2})$. Al dividir por $\rho=1+O(j^{-2})$ resulta
$\nu(c\mid B,\text{fila})=p_c(1+O(1/j))$ sobre todo el soporte de esa fila.

\paragraph{Lema de márgenes en el extremo (implicación revisada).}
Usamos la comparación de la ronda 10, parte 04, en el mismo soporte. Para constantes
$C,\gamma$ fijas, sean $\epsilon_T>0$ y $K_T$ su umbral y su constante relativa:
\[
 \left|P_e/p_e-1\right|\le K_T(\epsilon+1/j),\qquad
 \epsilon=\max_f\sum_{e\in f}p_e\le\epsilon_T .
\]
Aquí $P$ es el punto estacionario de torres; fuera del soporte ambas tablas valen cero.
Se pueden reservar de antemano $M_0=2C$ y $\gamma_T=\min(1,\gamma)/2$.
Si $L_{\max}<\lceil\alpha j\rceil$, por integralidad $L_{\max}<\alpha j$.
Así $\epsilon\le C\alpha$. Tomamos $C\alpha<\epsilon_T$, $K_TC\alpha\le1/4$
y después $K_T/j\le1/4$. Sobre el soporte $p_e/2\le P_e\le3p_e/2$. En consecuencia,
\[
 j\max P_e\le M_0,\qquad
 D\operatorname{diag}(P)D^t\succeq\tfrac12D\operatorname{diag}(p)D^t ,
 \qquad \max_f\sum_{e\in f}P_e\le M_0\alpha .
\]
$D$ es la incidencia de filas y columnas. Como $P$ es bistocástico, la matriz de
la izquierda tiene bloques diagonales identidad y cruzados $P,P^t$. Su espectro es
$1\pm\sigma_i(P)$; en el complemento del modo $(\mathbf1,-\mathbf1)$ se obtiene
$\|P-\mathbf1\mathbf1^t/j\|_{\rm op}\le1-\gamma_T$.

Para los menores hace falta el escalado bistocástico exacto de la máscara binaria $A$.
En el soporte, la estacionariedad de torres dice $P_{rc}(1-P_{rc})=a_r b_c$ con
factores positivos. Aplicando el lema de escalado de C, \S7, a $P$, su matriz
$P\circ(1-P)$ se escala por potenciales $O(1/j)$ en supremo. El resultado $D_A$
es un escalado de $A$. En el soporte $D_{A,e}/P_e=1+O(1/j)$ y fuera coinciden sus ceros.
La cota en Hilbert--Schmidt del mismo lema da $\|D_A-P\|_{\rm op}=O(1/j)$.
Se pueden reservar también $j\max D_{A,e}\le4C$ y
$\|D_A-\mathbf1\mathbf1^t/j\|_{\rm op}\le1-\gamma_T/2$ para $j$ grande.
En la llegada larga el uso de ese lema en el coste exterior emplea $p$, no $P$.

\paragraph{Uso del coste sobre una llegada regular larga.}
El enunciado de C, Corolario 6.8, dice primera salida. Para otro tablero fijo que cumple
las mismas cotas exteriores se usa su prueba. Escribimos $\ell_f=\sum_{e\in f}p_e$,
$h_f=1-\ell_f$, $X=Qe^{-L}$ y $\ell_f\le\varepsilon_{\rm out}$. La parte común en
el logaritmo del factor $F$ de C cumple, usando $\sum_f\ell_f=2j$,
\[
 \sum_f\log h_f-2\sum_e\log(1-p_e)
 \le 2\sum_e[-\log(1-p_e)-p_e]=O(1).
\]
En una permutación legal los $2j$ factores individuales cuestan a lo sumo
$\beta j$ con $\beta=2\varepsilon_{\rm out}/(1-\varepsilon_{\rm out})$.
El punto de reinas cumple difusión y brecha; es la tabla que usa C para los pesos
$W=p(1-p)$ de su permanente auxiliar. La estimación uniforme de ese permanente y
el cociente de normalizadores de C dan $X\le\exp(C_0+\beta j)$ si $Q>0$.
Para $Q=0$ se tiene $X=0$ directamente. La prueba no requiere el punto de torres
no ponderado de la máscara larga. Tampoco cambia $c>2\beta$ para las salidas históricas.
